% ============================================================
% Cheat sheet, argomento: il pattern Decorator (dispensa, capitolo 15)
% ============================================================
\argomento{Decorator}{Dispensa: capitolo 15. Pattern strutturale della GoF}[Pattern Decorator]

\scheda{Carta d'identità}{
  \pintento{Aggiungere \textbf{dinamicamente} responsabilità a un oggetto, \textbf{avvolgendolo} in un altro oggetto con
    la \textbf{stessa interfaccia}: un'alternativa flessibile alle sottoclassi.}
  \ptesto{«aggiungere funzionalità (ingredienti, optional, opzioni, filtri) a un oggetto», «a tempo di esecuzione»,
    «in qualunque combinazione, anche ripetuta» (la doppia mozzarella), «senza una sottoclasse per ogni
    combinazione» (con 10 ingredienti sarebbero $2^{10}$ classi), «un flusso di dati con buffer, cifratura,
    compressione».}
  \pcodice{una classe (spesso astratta) che \textbf{\code{implements} il Component} \textbf{e} ha \textbf{un} campo dello
    \textbf{stesso tipo}, ricevuto nel costruttore (\code{protected Pizza pizza}); i decoratori concreti chiamano il
    metodo dell'oggetto avvolto e \textbf{aggiungono} qualcosa (\code{pizza.getCost() + 1.50}); il client scrive
    catene come \code{new Mozzarella(new Fungo(new OlioEPomodoro()))}.}
  \pdiagramma{\code{Decorator} realizza (o estende) \code{Component} \textbf{e} ha un'aggregazione verso
    \code{Component} con molteplicità \code{1}; i \code{ConcreteDecorator} estendono il Decorator; il
    \code{ConcreteComponent} realizza il Component.}
  \pvarianti{il Decorator come classe astratta che delega tutto per default (i decoratori concreti ridefiniscono solo
    ciò che cambiano); nella libreria di Java i flussi di I/O: \code{new BufferedReader(new InputStreamReader(System.in))}
    (byte $\to$ caratteri $\to$ righe).}
  \pconfondere{il \textbf{Composite}: contiene \textbf{molti} componenti, il Decorator \textbf{uno}; l'\textbf{Adapter}:
    l'oggetto avvolto ha un'interfaccia \textbf{diversa} e l'adapter la traduce, il Decorator ha la \textbf{stessa}
    interfaccia e aggiunge; l'ereditarietà: aggiunge funzioni in modo statico, nel codice, non a tempo di esecuzione.}
}

\begin{struttura}{decorator-struttura}{La struttura}
\begin{tikzpicture}
  \umlclass[2.6cm]{c}{(1.6,0)}{\textit{Component}}{}{+ operation()}
  \umlclass[3.1cm]{cc}{(-0.6,-2.4)}{ConcreteComponent}{}{+ operation()}
  \umlclass[3.3cm]{d}{(3.4,-2.4)}{\textit{Decorator}}{- component: Component}{+ operation()}
  \umlclass[3.1cm]{cd}{(3.4,-4.9)}{ConcreteDecorator}{}{+ operation()}
  \draw[gen] (cc.north) |- ++(0,0.5) -| ([xshift=-0.45cm]c.south);
  \draw[gen] (d.north) |- ++(0,0.35) -| ([xshift=0.45cm]c.south);
  \draw[aggr] (d.east) -- ++(0.5,0) |- (c.east);
  \node[molt, anchor=south west] at ($(c.east)+(0.1,0.05)$) {1};
  \node[ruolo, anchor=north west] at ($(c.east)+(0.1,-0.05)$) {component};
  \draw[gen] (cd.north) -- (d.south);
  \node[notauml, anchor=east, text width=3.0cm] (n) at (1.1,-4.9) {component.operation();\\// più il comportamento\\// aggiunto};
  \draw[tratt] (cd.west) -- (n.east);
\end{tikzpicture}
\tcblower
\textbf{Component}: l'interfaccia comune. \textbf{ConcreteComponent}: l'oggetto di base. \textbf{Decorator}:
\textbf{è} un Component (eredita l'interfaccia) e \textbf{ha} un Component (il campo a cui delega): è questa doppia
relazione che permette di avvolgere un decoratore dentro un altro. \textbf{ConcreteDecorator}: chiama l'oggetto
avvolto e aggiunge il suo pezzo. La chiamata scende lungo la catena fino all'oggetto di base, e i risultati
risalgono.
\end{struttura}

\begin{programma}{decorator-esempio}{L'esempio del corso: la pizza}
\begin{lstlisting}[style=javacs, name=Pizzeria.java]
public class Pizzeria {
    public static void main(String[] args) {
        Pizza p = new OlioEPomodoro();
        stampa(p);
        p = new Mozzarella(p); // ogni new avvolge
        p = new Fungo(p);      // la pizza di prima
        p = new Mozzarella(p); // doppia mozzarella
        stampa(p);
    }
    static void stampa(Pizza p) {
        System.out.println(p.getDescription());
        System.out.println("  " + p.getCost());
    }
}
interface Pizza {                     // Component
    String getDescription();
    double getCost();
}
class OlioEPomodoro implements Pizza {   // di base
    public String getDescription() {
        return "Olio e Pomodoro";
    }
    public double getCost() { return 3.00; }
}
abstract class PizzaDecorator implements Pizza {
    protected final Pizza pizza;         // avvolta
    PizzaDecorator(Pizza pizza) {
        this.pizza = pizza;
    }
}
class Mozzarella extends PizzaDecorator {
    Mozzarella(Pizza p) { super(p); }
    public String getDescription() {
        return pizza.getDescription()
                + ", Mozzarella";
    }
    public double getCost() {
        return pizza.getCost() + 1.50;
    }
}
class Prosciutto extends PizzaDecorator {
    Prosciutto(Pizza p) { super(p); }
    public String getDescription() {
        return pizza.getDescription()
                + ", Prosciutto";
    }
    public double getCost() {
        return pizza.getCost() + 1.00;
    }
}
class Fungo extends PizzaDecorator {
    Fungo(Pizza p) { super(p); }
    public String getDescription() {
        return pizza.getDescription()
                + ", Funghi";
    }
    public double getCost() {
        return pizza.getCost() + 2.00;
    }
}
\end{lstlisting}
\tcblower
\begin{tikzpicture}
  \umlclass[3.6cm]{pz}{(1.2,0)}{\stereotipo{interface}{Pizza}}{}{+ getDescription(): String\\+ getCost(): double}
  \umlclass[3.4cm]{b}{(-0.8,-2.6)}{OlioEPomodoro}{}{+ getDescription(): String\\+ getCost(): double}
  \umlclass[3.8cm]{d}{(3.2,-2.6)}{\textit{PizzaDecorator}}{\# pizza: Pizza \{readOnly\}}{\textasciitilde\ PizzaDecorator(pizza: Pizza)}
  \node[classe1=1.7cm] (m) at (1.5,-4.7) {Mozzarella};
  \node[classe1=1.7cm] (r) at (3.3,-4.7) {Prosciutto};
  \node[classe1=1.4cm] (f) at (5.0,-4.7) {Fungo};
  \draw[realizza] (b.north) |- ++(0,0.5) -| ([xshift=-0.45cm]pz.south);
  \draw[realizza] (d.north) |- ++(0,0.35) -| ([xshift=0.45cm]pz.south);
  \draw[aggr] (d.east) -- ++(0.55,0) |- (pz.east);
  \node[molt, anchor=south west] at ($(pz.east)+(0.1,0.05)$) {1};
  \node[ruolo, anchor=north west] at ($(pz.east)+(0.1,-0.05)$) {pizza};
  \foreach \x in {m, r, f}{ \draw[gen] (\x.north) |- ($(d.south)+(0,-0.35)$) -- (d.south); }
\end{tikzpicture}

\medskip
\begin{lstlisting}[style=shellcs, name=Terminale]
$ javac Pizzeria.java
$ java Pizzeria
Olio e Pomodoro
  3.0
Olio e Pomodoro, Mozzarella, Funghi, Mozzarella
  8.0
\end{lstlisting}
\regola{\code{getCost()} sull'anello più esterno scende fino alla pizza base, che risponde 3,00; poi i risultati
risalgono e ogni decoratore aggiunge il suo pezzo: $3 + 1{,}5 + 2 + 1{,}5 = 8$. Le combinazioni (anche ripetute)
si scelgono a tempo di esecuzione con una classe per ingrediente. Il prezzo da pagare: ogni chiamata attraversa
tutta la catena.}
\end{programma}
